\documentclass[11pt,a4paper]{article}

% ── Packages ──────────────────────────────────────────────────
\usepackage[utf8]{inputenc}
\usepackage[T1]{fontenc}
\usepackage{amsmath,amssymb}
\usepackage{graphicx}
\usepackage[margin=2.3cm]{geometry}
\usepackage{booktabs}
\usepackage{caption}
\usepackage{subcaption}
\usepackage{hyperref}
\usepackage{float}
\usepackage{parskip}
\usepackage{enumitem}

\hypersetup{colorlinks=true, linkcolor=blue, citecolor=blue, urlcolor=blue}
\graphicspath{{plots/}}

\title{\textbf{Assignment 3: ARMA / Seasonal Processes and ARX Transfer-Function Models}\\[4pt]
       \large 02417 Time Series Analysis --- Spring 2026 \\
       \large Technical University of Denmark}
\author{%
  Jakub Piotrowski (s253074) \and
  Mohamed Amine Cheikh (s252221) \and
  Utkarsh Rupesh Saner (s252680)%
}
\date{\today}

\begin{document}
\maketitle

%===================================================================
\section{Stability of AR(2) processes}
%===================================================================
Let $\{X_t\}$ follow
\[
  X_t+\phi_1 X_{t-1}+\phi_2 X_{t-2}=\varepsilon_t,
  \qquad \varepsilon_t\sim\text{WN}(0,1).
\]
The characteristic polynomial is $\phi(z)=1+\phi_1 z+\phi_2 z^2$. The process
is stationary iff both roots of $\phi(z)=0$ lie outside the unit circle, which
in this sign convention is equivalent to the Schur conditions
\begin{equation}\label{eq:schur}
  |\phi_2|<1, \qquad \phi_1+\phi_2>-1, \qquad \phi_1-\phi_2<1.
\end{equation}
When stationary the theoretical ACF follows the Yule--Walker recursion
$\rho(1)=-\phi_1/(1+\phi_2)$ and $\rho(k)=-\phi_1\rho(k-1)-\phi_2\rho(k-2)$
for $k\ge 2$. For each parameter pair below we simulate 5 realisations of
length $n=200$ with a common seed and a 200-step burn-in, then compare the
empirical ACFs to $\rho(k)$ up to lag 30.

% -------- 1.1 / 1.2 --------
\subsection*{1.1--1.2~~$\phi_1=-0.6,\ \phi_2=0.5$}
Roots of $\phi(z)$: $0.6\pm 1.28i$, $|z|=\sqrt{2}\approx 1.414>1$. All three
Schur conditions in \eqref{eq:schur} hold, so the process is
\emph{stationary}. Because the discriminant $\phi_1^{2}-4\phi_2=-1.64<0$ the
roots are complex and the ACF exhibits damped oscillations.
\begin{figure}[H]\centering
  \begin{subfigure}{.48\linewidth}\centering\includegraphics[width=\linewidth]{p1_1_1_sims.pdf}\end{subfigure}\hfill
  \begin{subfigure}{.48\linewidth}\centering\includegraphics[width=\linewidth]{p1_1_1_acf.pdf}\end{subfigure}
  \caption{$\phi_1=-0.6,\phi_2=0.5$: 5 realisations (left) and empirical ACFs with the theoretical $\rho(k)$ overlaid (right).}
  \label{fig:sec1_1}
\end{figure}
The empirical ACFs cluster around $\rho(k)$ and reproduce the oscillating
decay with period $\approx 6$, matching the complex-root prediction.

% -------- 1.3 --------
\subsection*{1.3~~$\phi_1=-0.6,\ \phi_2=-0.3$}
Real roots $z\in\{1.08,-3.08\}$ are both outside the unit circle:
\emph{stationary}. The discriminant $\phi_1^{2}-4\phi_2=1.56>0$ so the ACF
decays monotonically. Empirical $\hat\rho(1)\approx 0.82$ is close to
$-\phi_1/(1+\phi_2)=0.857$.
\begin{figure}[H]\centering
  \begin{subfigure}{.48\linewidth}\centering\includegraphics[width=\linewidth]{p1_1_3_sims.pdf}\end{subfigure}\hfill
  \begin{subfigure}{.48\linewidth}\centering\includegraphics[width=\linewidth]{p1_1_3_acf.pdf}\end{subfigure}
  \caption{$\phi_1=-0.6,\phi_2=-0.3$: real-root stationary case with monotone ACF decay.}
\end{figure}

% -------- 1.4 --------
\subsection*{1.4~~$\phi_1=0.6,\ \phi_2=-0.3$}
Only the sign of $\phi_1$ is flipped compared to 1.3. Roots mirror to
$\{-1.08,\,3.08\}$, still stationary, and the ACF becomes monotone but
\emph{alternating in sign} because $\rho(1)=-\phi_1/(1+\phi_2)<0$
($\hat\rho(1)\approx -0.79$, theoretical $-0.857$). Visually the series
oscillates rapidly across its mean.
\begin{figure}[H]\centering
  \begin{subfigure}{.48\linewidth}\centering\includegraphics[width=\linewidth]{p1_1_4_sims.pdf}\end{subfigure}\hfill
  \begin{subfigure}{.48\linewidth}\centering\includegraphics[width=\linewidth]{p1_1_4_acf.pdf}\end{subfigure}
  \caption{$\phi_1=0.6,\phi_2=-0.3$: stationary but sign-alternating ACF.}
\end{figure}

% -------- 1.5 --------
\subsection*{1.5~~$\phi_1=-0.7,\ \phi_2=-0.3$}
Here $\phi_1+\phi_2=-1$, so the second Schur condition holds only on the
\emph{boundary}: one root equals exactly $1$ and the other $-10/3$. The
process is \emph{non-stationary} with a unit root. Simulated realisations
wander without returning to a mean (random-walk-like), the variance grows
roughly linearly with $t$, and the empirical ACF decays extremely slowly
($\hat\rho(10)\approx 0.70$). No theoretical $\rho(k)$ is plotted because it
is undefined.
\begin{figure}[H]\centering
  \begin{subfigure}{.48\linewidth}\centering\includegraphics[width=\linewidth]{p1_1_5_sims.pdf}\end{subfigure}\hfill
  \begin{subfigure}{.48\linewidth}\centering\includegraphics[width=\linewidth]{p1_1_5_acf.pdf}\end{subfigure}
  \caption{$\phi_1=-0.7,\phi_2=-0.3$: unit-root non-stationary process.}
\end{figure}

% -------- 1.6 --------
\subsection*{1.6~~$\phi_1=-0.75,\ \phi_2=-0.3$}
Now $\phi_1+\phi_2=-1.05<-1$, so the second Schur condition \emph{fails}. One
characteristic root has $|z|\approx 0.96<1$, placing it inside the unit
circle: the process is \emph{explosive}. The simulated sample variance grows
by several orders of magnitude over only 200 observations (measured mean
$\approx 4.7\times 10^{12}$). The empirical ACF again looks slow-decaying but
is meaningless for a non-stationary process.
\begin{figure}[H]\centering
  \begin{subfigure}{.48\linewidth}\centering\includegraphics[width=\linewidth]{p1_1_6_sims.pdf}\end{subfigure}\hfill
  \begin{subfigure}{.48\linewidth}\centering\includegraphics[width=\linewidth]{p1_1_6_acf.pdf}\end{subfigure}
  \caption{$\phi_1=-0.75,\phi_2=-0.3$: explosive non-stationary process.}
\end{figure}

% -------- 1.7 --------
\subsection*{1.7~~Time-series plot vs.\ ACF only}
Cases 1.5 and 1.6 show that the empirical ACF alone is a poor diagnostic for
stationarity: a slowly-decaying ACF is consistent with a stationary AR(2)
with real roots close to the unit circle (as in 1.3), with a unit-root
process (1.5), \emph{and} with a mildly explosive one (1.6). The raw
time-series plot, by contrast, immediately reveals drifting or exploding
behaviour (a visibly growing variance in 1.6 could never be mistaken for a
stationary process). One should therefore \emph{always} inspect both the
series itself and its ACF before drawing conclusions.

%===================================================================
\section{Predicting monthly solar power}
%===================================================================
The given model for the monthly generation $Y_t$ (MWh) is
\[
  (1+\phi_1 B)(1+\Phi_1 B^{12})\bigl(\log Y_t-\mu\bigr)=\varepsilon_t,
  \quad\phi_1=-0.38,\ \Phi_1=-0.94,\ \mu=5.72,\ \sigma_\varepsilon=0.22,
\]
with 36 monthly observations (2008-01 to 2010-12).
Introducing $X_t=\log Y_t-\mu$ and expanding the product of backshift factors,
\[
  (1+\phi_1 B)(1+\Phi_1 B^{12}) = 1+\phi_1 B+\Phi_1 B^{12}+\phi_1\Phi_1 B^{13},
\]
yields
\begin{equation}\label{eq:expand}
  X_t+\phi_1 X_{t-1}+\Phi_1 X_{t-12}+\phi_1\Phi_1 X_{t-13}=\varepsilon_t.
\end{equation}

% -------- 2.1 --------
\subsection*{2.1~~Residuals and model validation}
From \eqref{eq:expand} the one-step-ahead residuals are
\[
  \hat\varepsilon_{t+1|t}=X_{t+1}+\phi_1 X_{t}+\Phi_1 X_{t-11}+\phi_1\Phi_1 X_{t-12},
\]
which can be computed for $t=13,\dots,35$ (i.e.\ 23 residuals with indices
$14,\dots,36$). Their sample mean is $0.013$ and standard deviation $0.210$,
consistent with the assumed $\sigma_\varepsilon=0.22$. A battery of i.i.d.
tests yields

\begin{center}\small
\begin{tabular}{@{}lcccc@{}}\toprule
Ljung--Box $Q(5)$ & Ljung--Box $Q(10)$ & Shapiro--Wilk $W$ & max $|\hat\rho|$ (lags 1--10) & sample mean\\\midrule
$4.82\ (p=0.44)$ & $12.55\ (p=0.25)$ & $0.95\ (p=0.15)$ & $0.33$ (within 95\,\% band) & $0.013$ \\\bottomrule
\end{tabular}
\end{center}

No ACF coefficient exceeds the 95\,\% band, the Ljung--Box tests cannot
reject white noise at either horizon, and Shapiro--Wilk cannot reject
normality. The i.i.d. $\mathcal N(0,\sigma_\varepsilon^{2})$ assumption is
therefore compatible with the data (Fig.~\ref{fig:p2diag}). With only $n=23$
residuals the statistical power of all three tests is limited, so these
diagnostics are \emph{consistent with} --- rather than conclusive proof of
--- i.i.d. normality.

\begin{figure}[H]\centering
\includegraphics[width=.85\linewidth]{p2_residual_diagnostics.pdf}
\caption{Residual diagnostics for Model~\eqref{eq:expand}: time series, ACF, QQ-plot, and histogram.}
\label{fig:p2diag}
\end{figure}

% -------- 2.2 --------
\subsection*{2.2~~Twelve-month forecast}
Setting future innovations to their conditional expectation of zero in
\eqref{eq:expand} gives the recursion
\[
  \hat X_{t+k|t}=-\phi_1\,\hat X_{t+k-1|t}-\Phi_1\,\hat X_{t+k-12|t}-\phi_1\Phi_1\,\hat X_{t+k-13|t},
\]
with $\hat X_{t+j|t}=X_{t+j}$ for $j\le 0$. For $k\le 12$ both seasonal terms
use \emph{observed} values ($t+k-12\le t$ and $t+k-13\le t-1$), so only the
AR(1) component is genuinely recursive. Back-transformation uses
$\hat Y_{t+k|t}=\exp(\hat X_{t+k|t}+\mu)$ (no log-normal bias correction, as is
standard in this course). Results are in Table~\ref{tab:fc} and
Fig.~\ref{fig:p2fc}.

% -------- 2.3 --------
\subsection*{2.3~~95\,\% prediction intervals from the AR(1) part}
As instructed, the intervals rely on the AR(1) component only:
$X_t=a\,X_{t-1}+\varepsilon_t$ with $a=-\phi_1=0.38$. Its MA($\infty$)
representation has $\psi_j=a^j$ and equations (5.149)--(5.151) of the course
textbook give
\begin{equation}\label{eq:var}
  \operatorname{Var}\bigl(e_{t+k|t}\bigr)=\sigma_\varepsilon^{2}\sum_{j=0}^{k-1}\psi_j^{2}
    =\sigma_\varepsilon^{2}\,\frac{1-a^{2k}}{1-a^{2}}.
\end{equation}
A 95\,\% interval for $X_{t+k}$ is
$\hat X_{t+k|t}\pm 1.96\sqrt{\operatorname{Var}(e_{t+k|t})}$; applying
$\exp(\cdot+\mu)$ to both endpoints transforms it into an interval for
$Y_{t+k}$.

\begin{table}[H]\centering\small
\begin{tabular}{rlrrrr}\toprule
$k$ & month & $\hat Y_{t+k|t}$ [MWh] & 95\,\% lower & 95\,\% upper & $\sqrt{\operatorname{Var}(e_{t+k|t})}$\\\midrule
 1 & 2011-01 & 101.3 &  65.8 & 155.9 & 0.220\\
 2 & 2011-02 & 147.2 &  92.8 & 233.5 & 0.235\\
 3 & 2011-03 & 156.9 &  98.5 & 249.9 & 0.238\\
 4 & 2011-04 & 142.6 &  89.4 & 227.2 & 0.238\\
 5 & 2011-05 & 355.3 & 222.9 & 566.3 & 0.238\\
 6 & 2011-06 & 538.0 & 337.6 & 857.5 & 0.238\\
 7 & 2011-07 & 611.2 & 383.5 & 974.2 & 0.238\\
 8 & 2011-08 & 631.9 & 396.5 & 1007.2& 0.238\\
 9 & 2011-09 & 630.1 & 395.3 & 1004.3& 0.238\\
10 & 2011-10 & 508.2 & 318.8 & 809.9 & 0.238\\
11 & 2011-11 & 375.0 & 235.3 & 597.7 & 0.238\\
12 & 2011-12 & 280.5 & 176.0 & 447.1 & 0.238\\\bottomrule
\end{tabular}
\caption{Twelve-month forecast of monthly solar generation with 95\,\% PI derived from the AR(1) part of the model. The last column is the log-scale standard deviation $\sqrt{\operatorname{Var}(e_{t+k\mid t})}$ of the forecast error for $X_{t+k}=\log Y_{t+k}-\mu$; it saturates at $\sigma_\varepsilon/\sqrt{1-a^{2}}\approx 0.238$ by $k=3$.}\label{tab:fc}
\end{table}

\begin{figure}[H]\centering
\includegraphics[width=.85\linewidth]{p2_forecast.pdf}
\caption{Observed monthly generation (2008--2010) and 12-step-ahead forecast with 95\,\% PI (back-transformed).}\label{fig:p2fc}
\end{figure}

% -------- 2.4 --------
\subsection*{2.4~~Comment on forecast trustworthiness}
\textbf{Point forecast.} The predicted trajectory reproduces the expected
seasonal shape (low winter, August peak) and its level is in line with
2008--2010, so the point forecast is credible.

\textbf{Interval width.} The PI is too optimistic for at least three
reasons:
\begin{enumerate}[leftmargin=*,itemsep=0pt,topsep=2pt]
\item The AR(1) variance in \eqref{eq:var} \emph{saturates} almost immediately:
      $\sqrt{(1-a^{2k})/(1-a^{2})}\to 1/\sqrt{1-a^{2}}\approx 1.081$ for
      $k\ge 3$. This is why column~6 of Table~\ref{tab:fc} is essentially
      constant at $0.238$ from $k=3$ onward, whereas the true forecast
      uncertainty should keep growing.
\item The derivation \emph{ignores the seasonal $\Phi_1$ component}. For
      $k\le 12$ this is actually harmless --- all seasonal regressors
      $X_{t+k-12}$, $X_{t+k-13}$ draw from observed data and contribute no
      forecast noise, as we established in §2.2. But for $k\ge 13$ the
      $\Phi_1$-driven error term becomes the dominant source of uncertainty,
      and at that horizon the AR(1) approximation grossly underestimates
      the interval.
\item The three parameters $(\phi_1,\Phi_1,\mu)$ are \emph{treated as known}.
      With only 36 observations this is a strong assumption; parameter
      uncertainty would widen the real intervals further and make them grow
      with $k$.
\end{enumerate}
In short: the point forecast is trustworthy, the reported PI is a useful
lower bound on the actual uncertainty but is almost certainly too narrow and
the saturation behaviour is an artifact of using only the AR(1) part.

%===================================================================
\section{ARX model for the heating of a box}
%===================================================================
The data (231 hourly observations, file \texttt{box\_data\_60min.csv})
contain the heating power $P_{h,t}$ [W], the indoor--outdoor temperature
difference $T_{\Delta,t}$ [$^\circ$C] and the south-facing vertical solar
radiation $G_{v,t}$ [W/m$^2$]. Lags \texttt{.l0}--\texttt{.l10} are
pre-computed in the CSV.

% -------- 3.1 --------
\subsection*{3.1~~Exploratory time-series plot}
\begin{figure}[H]\centering
\includegraphics[width=.85\linewidth]{p3_timeseries.pdf}
\caption{Observed $P_h$, $T_\Delta$ and $G_v$ over the full 231-hour period.}
\end{figure}
$G_v$ shows a clear diurnal cycle peaking around noon. $T_\Delta$ tracks the
day--night cycle (warmer nights $\Rightarrow$ smaller $T_\Delta$). $P_h$ is
strongly anti-correlated with $G_v$ during daylight (incoming radiation heats
the box and the thermostat throttles the electrical heaters back) and
positively related to $T_\Delta$ (larger indoor--outdoor gradient requires
more heating). Visible strong persistence in all three series suggests
pronounced autocorrelation.

% -------- 3.2 --------
\subsection*{3.2~~Train/test split}
Rows $t=1,\dots,167$ (up to and including 2013-02-06 00:00) are used for
estimation; rows $t=168,\dots,231$ are held out as test set. Note that the
\texttt{thour} column in the CSV starts at 19 -- the ``$t_{\text{hour}}=1..167$''
notation in the assignment refers to the \emph{sequential row index} of the
training data, not to the \texttt{thour} value. The 1/64 normalisation used
later for the RMSE is consistent with the 64 test points.

% -------- 3.3 --------
\subsection*{3.3~~Investigation: scatter, ACF and CCF}
\begin{figure}[H]\centering
\includegraphics[width=.9\linewidth]{p3_eda.pdf}
\caption{Pairwise scatter plots (top), ACF of $P_h$ and CCFs of
$T_\Delta\to P_h$ and $G_v\to P_h$ (bottom). Training data only.}
\end{figure}
Scatter plots confirm the signs observed in 3.1: $P_h$ scales linearly and
positively with $T_\Delta$, and inversely (with some curvature) with $G_v$.
The ACF of $P_h$ decays slowly -- still $\approx 0.3$ at lag 20 -- reflecting
the thermal inertia of the box. The $T_\Delta\to P_h$ CCF is strongly positive
at small lags and slow-decaying; the $G_v\to P_h$ CCF is negative and
concentrated near lag 0, consistent with an almost instantaneous reduction
of heating by incoming radiation.

\emph{What cannot be read off directly from these plots}: (i) because the two
inputs are themselves autocorrelated (especially $G_v$ with its diurnal
cycle), the raw CCFs are contaminated by the inputs' own autocorrelation and
only a pre-whitening filter can isolate the true lag structure; (ii) the
CCFs cannot tell us whether a regression on contemporaneous values alone is
sufficient -- residual diagnostics are required for that (see~3.5).

\paragraph{Pre-whitened CCF (Box--Jenkins).}
To remove the confounding effect of the inputs' own autocorrelation we fit
an AR($p_x$) model to each input separately (order chosen by AIC over
$p\in\{1..10\}$) and apply the same filter
$\pi(B)=1-\hat\varphi_1 B-\cdots-\hat\varphi_{p_x} B^{p_x}$ to both the input
$x_t$ and to $P_{h,t}$. The CCF of the filtered series $\tilde u_t$,
$\tilde v_t$ in Fig.~\ref{fig:prewhiten} then has the standard
$\pm 1.96/\sqrt{n}$ band as its null distribution.
AIC selects $p_{T_\Delta}=2$ and $p_{G_v}=10$.
Both pre-whitened CCFs have their clearest peak at \emph{lag 0}
($|\hat\rho_{uv}(0)|\!\approx\!0.50$ for $T_\Delta$ and $0.85$ for $G_v$),
confirming that after removing the inputs' own dynamics, the dominant
transfer from either input to $P_h$ is essentially contemporaneous. This is
consistent with the raw CCF signs but is now statistically clean.
\begin{figure}[H]\centering
\includegraphics[width=.9\linewidth]{p3_prewhitened_ccf.pdf}
\caption{Pre-whitened CCFs from $T_\Delta$ and $G_v$ to $P_h$ after
AR-filtering each input (and $P_h$) with the AIC-selected filter of that
input. Dashed lines: $\pm 1.96/\sqrt{n}$.}\label{fig:prewhiten}
\end{figure}

% -------- 3.4 --------
\subsection*{3.4~~Impulse-response estimation}
Fitting a FIR regression
$P_{h,t}=c+\sum_{h=0}^{10}\beta_{h}^{T} T_{\Delta,t-h}+\sum_{h=0}^{10}\beta_{h}^{G} G_{v,t-h}+\varepsilon_t$
on the training set yields the impulse-response coefficients in
Fig.~\ref{fig:p3imp}.
\begin{figure}[H]\centering
\includegraphics[width=.85\linewidth]{p3_impulse.pdf}
\caption{Estimated impulse responses from $T_\Delta$ and $G_v$ to $P_h$ (lags 0--10).}
\label{fig:p3imp}
\end{figure}
The $T_\Delta\to P_h$ response has its largest weight at lag 1
($\approx 1.5$) and a secondary peak at lag 3, consistent with thermal
inertia (heat added now shows up in the heater set-point a step or two
later). The $G_v\to P_h$ response is largest in magnitude at lag 0
($\approx -0.10$), decays quickly, and is practically zero beyond lag 3:
solar gain acts on the heater almost contemporaneously.
A joint FIR (rather than two univariate pre-whitened impulse responses)
is used here because $T_\Delta$ and $G_v$ share a strong common diurnal
component; univariate pre-whitening would mis-attribute that common
variance to whichever input is processed first.

% -------- 3.5 --------
\subsection*{3.5~~Static regression}
The static model $P_{h,t}=\omega_1 T_{\Delta,t}+\omega_2 G_{v,t}+\varepsilon_t$
estimated on the training set gives
\[
  \hat\omega_1=+3.89\ (\mathrm{SE}\,0.023),\qquad
  \hat\omega_2=-0.110\ (\mathrm{SE}\,0.004),\qquad
  \hat\sigma=5.41,\qquad R^{2}=0.887.
\]
Both coefficients are massively significant ($|t|>26$).
\begin{figure}[H]\centering
\includegraphics[width=.9\linewidth]{p3_static_resid.pdf}
\caption{Static-regression residuals: time series, ACF, and CCFs with both inputs.}
\end{figure}
The signs match the physics (more heating with a larger thermal gradient,
less heating with incoming radiation) and 88.7\,\% of the variance is
explained. However the residual ACF is massively outside the 95\,\% band at
short lags and the residual CCFs with both inputs retain structure. The
static model fails to capture the box's thermal memory: a dynamic /
transfer-function (ARX) model is clearly required.

% -------- 3.6 --------
\subsection*{3.6~~First-order ARX}
Fitting
\[
  P_{h,t}=-\phi_1 P_{h,t-1}+\omega_1 T_{\Delta,t}+\omega_2 G_{v,t}+\varepsilon_t
\]
on the training set yields
\[
  \hat\phi_1=-0.419\ (\mathrm{SE}\,0.018),\quad
  \hat\omega_1=+2.31\ (\mathrm{SE}\,0.071),\quad
  \hat\omega_2=-0.084\ (\mathrm{SE}\,0.002),\quad \hat\sigma=2.67.
\]
All three coefficients are strongly significant.
The training RSS drops from $4883$ (static) to $1193$ (ARX(1)) --- a
$\approx 75\,\%$ reduction --- and the residuals look much closer to white
noise (Fig.~\ref{fig:arx1}). A large improvement is obtained from a single
AR lag alone.
\begin{figure}[H]\centering
\includegraphics[width=.9\linewidth]{p3_arx1_resid.pdf}
\caption{ARX(1) residuals: the residual ACF is mostly inside the 95\,\% band and the CCFs are visibly reduced compared to Section 3.5.}
\label{fig:arx1}
\end{figure}

% -------- 3.7 --------
\subsection*{3.7~~Model-order selection via AIC / BIC}
For $p=1,\dots,10$ we fit the ARX$(p)$ model
\[
  P_{h,t}=-\sum_{i=1}^{p}\phi_i P_{h,t-i}
          +\sum_{j=0}^{p-1}\omega_{1,j}T_{\Delta,t-j}
          +\sum_{j=0}^{p-1}\omega_{2,j}G_{v,t-j}+\varepsilon_t,
\]
i.e.\ the same number of lags for all variables (giving $k=3p$ free
parameters), and compute
$\mathrm{AIC}=n\log\hat\sigma^{2}+2k$ and
$\mathrm{BIC}=n\log\hat\sigma^{2}+k\log n$ with $n=167$.
\begin{figure}[H]\centering
\includegraphics[width=.95\linewidth]{p3_order_selection.pdf}
\caption{AIC and BIC (left) and one-step test-set RMSE (right) versus ARX order $p$.}\label{fig:ic}
\end{figure}
\begin{table}[H]\centering\small
\begin{tabular}{r r r r r r}\toprule
$p$ & $k$ & AIC & BIC & $\hat\sigma_{\text{train}}$ & RMSE$_{\text{test}}$ [W]\\\midrule
 1 &  3 & 334.3 & 343.7 & 2.67 & 3.75\\
 2 &  6 & 284.2 & 302.9 & 2.26 & 3.25\\
 3 &  9 & 265.6 & 293.6 & 2.10 & \textbf{2.88}\\
 4 & 12 & 246.0 & 283.4 & 1.94 & 3.04\\
 5 & 15 & 233.5 & \textbf{280.2} & 1.84 & 3.13\\
 6 & 18 & \textbf{227.3} & 283.5 & 1.77 & 3.33\\
 7 & 21 & 231.4 & 296.9 & 1.76 & 3.32\\
 8 & 24 & 234.6 & 309.4 & 1.75 & 3.27\\
 9 & 27 & 236.9 & 321.1 & 1.73 & 3.26\\
10 & 30 & 231.7 & 325.2 & 1.67 & 3.30\\\bottomrule
\end{tabular}
\caption{Information criteria and one-step test RMSE. Bold: column minima.}\label{tab:ic}
\end{table}
AIC, whose parameter penalty is only $+2k$, keeps decreasing until $p=6$.
BIC uses the heavier penalty $k\log n=k\log 167\approx 5.12\,k$ and picks
the simpler $p=5$. The two curves diverge from $p\approx 3$ onward because
the in-sample fit keeps improving but the BIC penalty grows faster than the
AIC penalty. Based on parsimony and the usual sample-size consistency
argument, BIC's $p=5$ is our preferred order.

% -------- 3.8 --------
\subsection*{3.8~~One-step predictions on the test set}
Using the coefficients fitted on the training set, we compute one-step
predictions on the test set (true lagged $P_h$ is used -- no simulation yet)
and evaluate
\[
  \mathrm{RMSE}=\sqrt{\tfrac{1}{64}\sum_{t=168}^{231}\hat\varepsilon_{t|t-1}^{\,2}}.
\]
The minimum of $\mathrm{RMSE}_{\text{test}}$ is attained at $p=3$
($=2.88$~W; see Table~\ref{tab:ic} and Fig.~\ref{fig:ic}); it then rises
monotonically, confirming over-fitting at higher orders. The honest
out-of-sample criterion agrees much better with BIC than with AIC: the
BIC-optimal $p=5$ is close to the RMSE optimum ($3.13$ vs.\ $2.88$~W),
whereas AIC's $p=6$ is clearly worse on the test set ($3.33$~W). We
therefore adopt $p^{\star}=5$ for the subsequent simulation.

% -------- 3.9 --------
\subsection*{3.9~~Multi-step simulation}
Before simulating we verify that the residuals of the final ARX($p^{\star}=5$)
model are compatible with white noise on the training set
(Fig.~\ref{fig:arx5}). The Ljung--Box statistics are
$Q(10)=6.46\ (p=0.78)$ and $Q(20)=13.44\ (p=0.86)$, both well above the
usual 0.05 threshold; the residual--input CCFs stay bounded by
$\max|\rho|\approx 0.20$ over $|h|\le 20$, i.e.\ essentially flat compared
to the raw $P_h$ CCFs in Fig.~\ref{fig:p3imp}. The identified model has
therefore captured both the autoregressive dynamics of $P_h$ and the transfer
from the two inputs; no further structure is left to exploit.
\begin{figure}[H]\centering
\includegraphics[width=.9\linewidth]{p3_arx5_resid.pdf}
\caption{Residual diagnostics for the final ARX(5) model: time trace, ACF, and CCFs with both inputs. Ljung--Box $Q(10)=6.46\,(p=0.78),\ Q(20)=13.44\,(p=0.86)$.}\label{fig:arx5}
\end{figure}

With $p^{\star}=5$ we simulate $P_h$ over the \emph{entire} period: the
observed inputs $T_\Delta,G_v$ are fed at each step, while the $P_h$ lags are
taken from previously simulated values (the first $p^{\star}$ values are
initialised from the observation).
\begin{figure}[H]\centering
\includegraphics[width=.9\linewidth]{p3_multistep_sim.pdf}
\caption{Observed (black) and multi-step simulated (red) $P_h$ for ARX(5). The dashed line marks the train/test split.}
\end{figure}
The simulation RMSE is $2.11$~W on the training period and $3.73$~W on the
test period, only marginally larger than the one-step RMSE at the same
order ($3.13$~W). The model reproduces the diurnal cycle and the overall
level across the full 231 hours without drifting, because the observed
inputs (driven by weather) continuously inject new information even while
the AR lags are allowed to propagate.

\textbf{Operational feasibility.} Such multi-step predictions \emph{are}
feasible in a real-time setting provided that forecasts of the inputs are
available. In this particular experiment $T_\Delta$ also depends on the
indoor temperature which -- being thermostat-controlled -- is essentially
constant, so effectively only external-temperature and solar-radiation
forecasts (standard outputs of any numerical weather prediction system) are
required.

% -------- 3.10 --------
\subsection*{3.10~~Conclusion}
\begin{enumerate}[leftmargin=*,itemsep=2pt,topsep=2pt]
\item The static regression already explains $88.7\,\%$ of the variance, but
      its residuals are strongly autocorrelated -- a dynamic /
      transfer-function model is needed.
\item Adding a single AR lag (ARX(1)) cuts the training RSS by $\approx 75\,\%$
      and largely whitens the residuals. Additional orders keep improving
      the in-sample fit, but AIC (p=6), BIC (p=5) and one-step test RMSE
      (p=3) give different optima; BIC is preferred here since its choice
      is closest to the out-of-sample optimum.
\item With $p^{\star}=5$ a multi-step simulation that uses only the observed
      inputs tracks $P_h$ well over the full 231 hours (train/test RMSE
      $2.11$ / $3.73$ W), showing that the identified model is usable in
      real-time predictive control given forecasts of $T_\Delta$ and $G_v$.
\item \emph{Limitations.} Effective solar gain depends nonlinearly on sun
      angle and surface orientation (not modelled here); the calibration
      period is short ($\approx 7$ days); and no uncertainty band is
      attached to the multi-step predictions. Natural extensions would be
      (i) an ARMAX formulation giving analytic $k$-step variance, or
      (ii) bootstrapping the training residuals to obtain an empirical
      prediction band around the simulation in Fig.~\ref{fig:arx5}.
\end{enumerate}

\end{document}
